\subsection{Central Functions on a Finite Group}

PROPOSITION 5. --- The family \((\chi_{\lambda})_{\lambda\in\widehat{\mathbf{G}}}\) is a basis of the vector space of central functions. The number of classes of simple linear representations of G is equal to the number of conjugacy classes of G.

For \(a = \sum_{g \in G} a_g g \in K[G]\) , write \(a^\vee = \sum_{g \in G} a_{g^{-1}} g\) ; the mapping \(a \mapsto a^\vee\) is an involutive antiautomorphism of the algebra K{[}G{]}. By formula (17), we have \(e_\lambda = |G|^{-1} d_\lambda \chi_\lambda^\vee\) for every \(\lambda \in \widehat{G}\) . The proposition then follows from the fact that the family \((e_\lambda)_{\lambda \in \widehat{G}}\) is a basis of the center of K{[}G{]}.

Denote by C the set of conjugacy classes of G. Let C be an element of C; the image \(C^{-1}\) of C by the mapping \(g \mapsto g^{-1}\) is a conjugacy class. The commutants of the elements of C are mutually conjugate subgroups of G; their cardinal \(d(\mathrm{C})\) satisfies

\[
\operatorname{Card} (\mathrm{G}) = \operatorname{Card} (\mathrm{C}) d (\mathrm{C}).\tag{30}
\]

In particular, we have \(d(\mathrm{C}) \cdot 1 \neq 0\) in the field K.

Let f be a central function on G. For any conjugacy class C, denote by \(f(\mathrm{C})\) the common value of the \(f(x)\) for \(x \in C\) . With this notation, the orthogonality relation for characters (VIII, p.~410, Proposition 4) can be written as

\[
\sum_ {\mathrm{C} \in \mathcal {C}} \chi_ {\lambda} (\mathrm{C} ^ {- 1}) \chi_ {\mu} (\mathrm{C}) d (\mathrm{C}) ^ {- 1} = \delta_ {\lambda \mu}\tag{31}
\]

for \(\lambda\) and \(\mu\) in \(\widehat{G}\) .

Denote by \(A\) the matrix of type \(\widehat{\mathbf{G}} \times \mathcal{C}\) with entries \(\chi_{\lambda}(\mathrm{C})\) and by \(B\) the matrix of type \(\mathcal{C} \times \widehat{\mathbf{G}}\) with entries \(\chi_{\lambda}(\mathrm{C}^{-1}) d(\mathrm{C})^{-1}\) . The sets \(\widehat{\mathbf{G}}\) and \(\mathcal{C}\) have the same cardinal (Proposition 5); relation (31) expresses the fact that the matrix product \(AB\) is the matrix unit of type \(\widehat{\mathbf{G}} \times \widehat{\mathbf{G}}\) . By Proposition 11 of II, §10, No.~12, p.~360, the matrix product \(BA\) is the matrix unit of type \(\mathcal{C} \times \mathcal{C}\) ; in other words, we have the relation

\[
\sum_ {\lambda \in \widehat {\mathrm{G}}} \chi_ {\lambda} (\mathrm{C} ^ {- 1}) \chi_ {\lambda} (\mathrm{C} ^ {\prime}) = d (\mathrm{C}) \delta_ {\mathrm{CC} ^ {\prime}}\tag{32}
\]

for C and C' in C (sometimes called the ``second orthogonality relation for characters'').

Let H be a subgroup of G. Let us note that the integer \(\text{Card(H)}\) divides \(\text{Card(G)}\) and that \(|G|\) is not zero in K; we therefore have \(|H| \neq 0\) in K.

Denote by \(Res_{H}^{G}\) the linear mapping from \(Z(K[G])\) to \(Z(K[H])\) that sends a central function on G to its restriction to H. We have seen that if \(\chi_{\pi}\) is the character of a finite-dimensional representation \(\pi\) of G, then \(\mathrm{Res}_{\mathrm{H}}^{\mathrm{G}}(\chi_{\pi})\) is the character of the representation \(\mathrm{Res}_{\mathrm{H}}^{\mathrm{G}}(\pi)\) of H.

PROPOSITION 6. --- Let f be a central function on G and u be a central function on H. We have

\[
\langle \mathrm{Ind} _ {\mathrm{H}} ^ {\mathrm{G}} (u), f \rangle_ {\mathrm{G}} = \langle u, \mathrm{Res} _ {\mathrm{H}} ^ {\mathrm{G}} (f) \rangle_ {\mathrm{H}}.\tag{33}
\]

The characters of the simple representations of G form a basis of \(Z(\mathrm{K}[\mathrm{G}])\) (VIII, p.~411, Proposition 5), and the same holds for H. It therefore suffices to establish (33) in the case when f is the character \(\chi_{\pi}\) of a simple representation \(\pi\) of G and u is the character \(\chi_{\sigma}\) of a simple representation \(\sigma\) of H. In this case, \(\operatorname{Ind}_{\mathrm{H}}^{\mathrm{G}}(u)\) is the character of the representation \(\operatorname{Ind}_{\mathrm{H}}^{\mathrm{G}}(\sigma)\) of G, and, by VIII, p.~410, Corollary, we have

\[
\langle \mathrm{Ind} _ {\mathrm{H}} ^ {\mathrm{G}} (u), f \rangle_ {\mathrm{G}} = (\dim_ {\mathrm{K}} \mathrm{Hom} _ {\mathrm{G}} (\mathrm{Ind} _ {\mathrm{H}} ^ {\mathrm{G}} (\sigma), \pi)) \cdot 1.
\]

We show the relation

\[
\langle u, \mathrm{Res} _ {\mathrm{H}} ^ {\mathrm{G}} (f) \rangle_ {\mathrm{H}} = (\dim_ {\mathrm{K}} \mathrm{Hom} _ {\mathrm{H}} (\sigma , \mathrm{Res} _ {\mathrm{H}} ^ {\mathrm{G}} (\pi))) \cdot 1
\]

likewise, and equality (33) follows by Frobenius reciprocity.

PROPOSITION 7. --- Let \(f\) be a mapping from \(G\) to \(K\) . The following assertions are equivalent:

\begin{enumerate}
\def\labelenumi{(\roman{enumi})}
\item
  There are an element \(\lambda\) of \(\widehat{\mathbf{G}}\) and an element \(a\) of \(\mathrm{K}\) such that \(f = a\chi_{\lambda}\) .
\item
  For every pair \((g,g^{\prime})\) of elements of \(\mathbf{G}\) , we have
\end{enumerate}

\[
f (g) f (g ^ {\prime}) = | \mathrm{G} | ^ {- 1} f (1) \sum_ {h \in \mathrm{G}} f (h g h ^ {- 1} g ^ {\prime}).\tag{34}
\]

Let \(\lambda\) be an element of \(\widehat{\mathbf{G}}\) ; for every endomorphism \(u\) of \(\mathrm{V}_{\lambda}\) , set

\[
u ^ {\natural} = | \mathrm{G} | ^ {- 1} \sum_ {h \in \mathrm{G}} \pi_ {\lambda} (h) u \pi_ {\lambda} (h ^ {- 1}).\tag{35}
\]

The endomorphism \(u^{\natural}\) of \(V_{\lambda}\) is K{[}G{]}-linear. By Schur's lemma (VIII, p.~47, Theorem 1), \(u^{\natural}\) is a homothety. Since \(u\) and \(u^{\natural}\) have the same trace, we therefore have

\[
u ^ {\natural} = d _ {\lambda} ^ {- 1} \mathrm{Tr} (u) 1 _ {\mathrm{V} _ {\lambda}}.
\]

Let \(u\) and \(v\) be endomorphisms of \(V_{\lambda}\) ; it follows that

\[
\operatorname{Tr} (u) \operatorname{Tr} (v) = d _ {\lambda} \operatorname{Tr} (u ^ {\natural} v) = d _ {\lambda} | G | ^ {- 1} \sum_ {h \in G} \operatorname{Tr} (\pi_ {\lambda} (h) u \pi_ {\lambda} (h ^ {- 1}) v).\tag{36}
\]

Take \(u = \pi_{\lambda}(g)\) and \(v = \pi_{\lambda}(g')\) in formula (36); relation (34) for \(f = \chi_{\lambda}\) follows.

Conversely, suppose that assertion (ii) holds. If we have \(f(1)=0\) , then it follows that \(f(g)f(g')=0\) for every pair \((g,g')\) of elements of G, and therefore f=0. We can therefore assume \(f(1)\neq0\) . Taking \(g'=1\) in (34), we obtain the relation

\[
f (g) = | \mathrm{G} | ^ {- 1} \sum_ {h \in \mathrm{G}} f (h g h ^ {- 1})
\]

for every \(g \in G\) , which implies that f is a central function. By Proposition 5, there exists a family \((a_{\lambda})\) of elements of K such that \(f = \sum_{\lambda \in \widehat{G}} a_{\lambda} \chi_{\lambda}\) . Let us replace f with this expression in formula (34); taking into account that each of the characters \(\chi_{\lambda}\) also satisfies this relation, we find

\[
\sum_ {\lambda , \mu \in \widehat {G}} a _ {\lambda}   a _ {\mu}   \chi_ {\lambda} (g)   \chi_ {\mu} (g ^ {\prime}) = \sum_ {\lambda \in \widehat {G}} a _ {\lambda}   d _ {\lambda} ^ {- 1}   f (1)   \chi_ {\lambda} (g)   \chi_ {\lambda} (g ^ {\prime})\tag{37}
\]

for \(g, g' \in \mathbf{G}\) . This relation can also be written as

\[
\sum_ {\lambda , \mu} (a _ {\lambda} a _ {\mu} - \delta_ {\lambda \mu} a _ {\lambda} d _ {\lambda} ^ {- 1} f (1)) \chi_ {\lambda} (g) \chi_ {\mu} (g ^ {\prime}) = 0\tag{38}
\]

for \(g, g' \in \mathrm{G}\) . Now, the functions \(\chi_{\lambda}\) , for \(\lambda \in \widehat{\mathrm{G}}\) , are linearly independent (Proposition 5 of VIII, p.~411); it follows that

\[
a _ {\lambda} a _ {\mu} = \delta_ {\lambda \mu} a _ {\lambda} d _ {\lambda} ^ {- 1} f (1)
\]

for \(\lambda, \mu \in \widehat{\mathbf{G}}\) . In particular, \(a_{\lambda}a_{\mu} = 0\) whenever \(\lambda \neq \mu\) . Consequently, there exists at most one element \(\lambda\) of \(\widehat{\mathbf{G}}\) such that \(a_{\lambda} \neq 0\) , and we have \(f = a_{\lambda}\chi_{\lambda}\) , and therefore (i).

